\section{Summary}
This thesis built a trimodal emotion recognition system on the EAV corpus---AST for speech, a facial-emotion ViT for video, EEGNet for EEG, and an attention-based fusion module over the three---and asked when learned, attention-based fusion earns its complexity over simply averaging the three modalities' predictions.

The answer depends on two things: how much data the fusion module has, and whether its input still encodes who the participant is. Within-subject, with 280 training trials per participant, learned fusion is at best marginally better than averaging once test-set epoch selection is removed (concat-MLP 79.96\% and modality-dropout attention 79.80\% against 77.48\%), and plain attention is not better at all. On unseen participants with default normalisation, every learned head is 3.4--4.5 points \emph{worse} than averaging (62.2\%). With a label-free, per-participant normalisation computed from about 20 unlabelled clips, every learned head is 11--14 points \emph{better} (concat-MLP 75.9\%, modality-dropout attention 75.0\%, attention 73.9\%), and linear probes show that what the normalisation removes is participant identity, which is decodable at 82--100\% from every modality beforehand and at chance afterwards.

The thesis also re-examined its own Pre-Thesis~II results. Selecting each fusion head's epoch on the test set had inflated them by 5.3--7.5 points; the claimed state-of-the-art figure is withdrawn. Trial correspondence across modalities cannot be verified from the released data. And the suppression matrix that Pre-Thesis~II presented as a map of emotion suppression is statistically indistinguishable from the EEG classifier's own error structure.

\section{Summary of Contributions}

\textbf{Contribution 1: Unsupervised subject calibration for trimodal fusion.} Per-participant standardisation of the fusion features, from an unlabelled recording of the new user, turns learned fusion from the worst to the best option on unseen participants on EAV (Chapter~\ref{chap:calibration}). The contribution includes the operating characteristics a deployment needs: significance for every head from 20 clips (100 seconds), saturation at about 50, failure of a neutral-only enrolment session, and the need for 50--100 clips when the session is emotionally unbalanced. It also includes the mechanism: identity is linearly decodable from every modality under the default normalisation and at chance after calibration, and vision---the modality whose emotion signal is most entangled with identity---gains most. Under calibration the fusion heads rank concat-MLP $>$ modality-dropout attention $>$ attention, each step significant.

\textbf{Contribution 2: An evaluation-integrity audit of trimodal fusion on EAV.} Three cheap tests---a same-trajectory comparison of test-selected and leak-free readouts, an element-wise label-sequence check that reports whether it can fail, and a class-conditional permutation null for coherence matrices---overturn or weaken the most prominent Pre-Thesis~II claims (Chapter~\ref{chap:audit}). The tests run in minutes on cached artefacts and apply directly to other work on this benchmark, none of which reports its selection protocol.

\textbf{Contribution 3: The first subject-independent evaluation of trimodal fusion on EAV.} Five-fold subject-wise cross-validation with every encoder retrained per fold shows that audio transfers to new participants at parity, EEG loses 5.7 points and vision loses 33.4 (Chapter~\ref{chap:crosssubject}).

\textbf{Contribution 4: A trimodal attention fusion module with a zero-EEG inference path.} \texttt{TrimodalAttentionFusion} exposes per-modality embeddings for future temporal heads and, trained with softhard modality dropout, remains accurate with EEG absent: within-subject, it loses 0.63 points without EEG and still beats attention trained without dropout (Chapter~\ref{chap:within}).

\textbf{Contribution 5: Documented fixes to the EAV reference code.} Defects in the public PyTorch EEGNet path---among them a forward hook that replaced a layer's output with its weight tensor, and a trainer that left BatchNorm and dropout in evaluation mode after the first epoch---prevented EEGNet from training correctly. With the defects fixed and the full epoch budget restored, pilot EEG accuracy rose from 24.4\% to 50.6\% (Chapter~\ref{chap:engineering_challenges}).

\section{Answers to the Research Questions}

\begin{table}[htbp]
\centering
\small
\caption{Summary answers to the nine research questions.}
\label{tab:rq_summary}
\begin{tabular}{p{2.6cm} p{11.0cm}}
\hline
\textbf{Question} & \textbf{Answer} \\
\hline
RQ1 Baselines & Yes: audio 57.1\%, vision 74.6\%, EEG 43.9\% within-subject, all above the EAV baselines. \\
RQ2 Fusion gain & Yes: naive late fusion 77.5\% beats vision 74.6\% ($p = 0.020$). \\
RQ3 Architecture & Within-subject, plain attention does not beat averaging; concat-MLP and dropout-trained attention beat it by about 2.4 points, marginally. \\
RQ4 Coherence & The patterns reproduce EEG's error structure ($r = 0.989$) and do not exceed a class-conditional permutation null within-subject. One cross-subject candidate (Sadness$\rightarrow$Neutral) remains. \\
RQ5 Robustness & Yes: without EEG the dropout model scores 79.17\% within-subject (not significantly below full) and 73.57\% cross-subject when calibrated. \\
RQ6 Generalisation & 62.2\% on unseen participants (chance 20\%); the loss is concentrated in vision. \\
RQ7 Protocol dependence & Learned fusion loses to averaging on unseen participants with default normalisation and wins by 11--14 points with calibration. \\
RQ8 Calibration & Yes, from about 20 unlabelled clips, provided the clips are emotionally varied; the mechanism is removal of participant identity. \\
RQ9 Integrity & Selection inflated fusion results by 5.3--7.5 points; alignment is unverified; the suppression claims do not survive. \\
\hline
\end{tabular}
\end{table}

\section{Limitations}
Seven limitations are acknowledged.

\textbf{Limitation 1: Controlled corpus, not spontaneous data.} EAV uses cue-based elicitation. Emotions are performed on instruction, the audio-visual consensus is usually the cued label, and the EEG and audio-visual windows are not simultaneous. All results describe this setting.

\textbf{Limitation 2: A single corpus and a single session.} All results, under every protocol, come from one laboratory corpus with one recording configuration. In the calibration study, calibration and evaluation clips come from the same session; whether calibration statistics transfer across days, rooms or devices is untested.

\textbf{Limitation 3: The calibrated comparison has an uncalibrated baseline.} Per-participant normalisation acts on the fusion features, which naive late fusion does not use. An analogous per-participant correction of the encoder logits was not evaluated, so the 11--14-point margin is between calibrated learned fusion and uncalibrated averaging. Calibrated single-modality linear probes reach at most 67.8\%, which suggests that most, but not necessarily all, of the margin is due to fusion.

\textbf{Limitation 4: Calibration requires enrolment.} The method needs a short unlabelled recording of each new user, with varied emotional content, before it can normalise their inputs. It does not support one-shot inference on a single clip from an unseen person.

\textbf{Limitation 5: Pooled-vector features.} The encoders produce one pooled vector per modality per trial, so the attention module relates only three tokens. Sequence-level features and pairwise cross-modal attention in the style of Lee, Kim and Kim \cite{Lee2024MERCL} were not evaluated; the finding that concat-MLP outperforms attention may not hold at finer temporal resolution.

\textbf{Limitation 6: Unverified trial alignment.} The released pickles are sorted into class blocks and carry no trial identifiers, so it cannot be confirmed that trial $i$ is the same interaction in every modality. Classification results do not depend on this; trial-level coherence analyses do.

\textbf{Limitation 7: Unverified EEGNet port.} The corrected PyTorch EEGNet was not compared against the TensorFlow implementation from which the EAV paper's EEG baseline most likely comes, so the comparison with that baseline is between implementations as well as between training configurations.

\section{Future Work}

\textbf{Phase 2.1: Close the calibration comparison.} The most direct open question is Limitation~3: applying a per-participant correction to the encoder logits before averaging, and comparing calibrated averaging with calibrated learned fusion on the same folds. The cached features, logits and trained heads needed for this already exist, so the experiment requires no retraining.

\textbf{Phase 2.2: Cross-session calibration.} Recording the same participants on two occasions, calibrating on one and evaluating on the other, would test whether the enrolment statistics describe the person or only the session.

\textbf{Phase 2.3: Coherence with an independent internal-state measure.} The one surviving coherence candidate, Sadness$\rightarrow$Neutral under subject-independent evaluation, can only be interpreted against a measure of what the participant felt. A study collecting per-trial self-report alongside EEG, analysed with the class-conditional permutation null of Chapter~\ref{chap:audit}, is the minimum design under which a suppression claim could be supported.

\textbf{Phase 2.4: Primary audio-visual dataset and merged training.} A primary recording study would collect audio and video from 5--10 new subjects performing the same five-emotion task as EAV, recorded with consumer-grade smartphones in non-laboratory settings. The primary dataset would be merged with EAV through two-stage fine-tuning: the AST and ViT encoders are pre-trained on EAV's audio and speaking-clip vision data, then fine-tuned on the union of EAV speaking clips and the new primary recordings with appropriate per-corpus weighting. The EEG branch would continue to train on EAV alone, with the primary dataset contributing only audio-visual examples to the fusion module. At inference on primary recordings, the EEG token is zeroed, matching the zero-EEG configuration of Section~\ref{sec:modality_dropout_results}. This design tests two properties not testable on EAV alone: cross-corpus generalisation of the audio-visual encoders, and the ecological validity of the zero-EEG demo path on data recorded outside the EAV laboratory protocol. Cross-modal coherence on the primary dataset would substitute self-reported internal-state Likert ratings, collected per trial, for the unavailable EEG signal; cross-modal disagreement is then measured between audio-vision consensus and self-report. EEG instrumentation, ethics approval for primary EEG recording, and behavioural-EEG validation of the Likert proxy are required for a full trimodal extension that includes a primary EEG channel.

\textbf{Phase 2.5: In-the-wild temporal emotion tracking.} The architecture designed in this thesis---with per-modality embeddings exposed as named outputs---supports attachment of a temporal head, such as an LSTM or Residual-TCN, for continuous emotion tracking on temporally annotated video corpora such as DFEW or MAFW. This extension requires a temporally annotated in-the-wild corpus, a voice activity detector for audio, a robust face detector and tracker for video, and a temporal decoder consuming the per-modality embedding sequences. The per-modality encoders (AST, ViT, EEGNet) can be frozen, making this an incremental addition rather than a full redesign.

\textbf{Phase 2.6: MERCL contrastive pre-training.} The supervised contrastive pre-training scheme of Lee, Kim and Kim \cite{Lee2024MERCL}---comprising intra-modal contrastive loss (AMCL), inter-modal alignment loss (EMCL), and sample-wise alignment (SMCL)---was dropped from the present thesis under schedule pressure. Its re-introduction in Phase 2 would test whether cross-modal contrastive alignment of the per-modality encoder representations, before fine-tuning the fusion head, materially improves classification accuracy, and in particular whether an identity-invariant contrastive objective can replace the enrolment session that subject calibration requires. The hypothesis, based on the Lee et al. ablation results, is that contrastive pre-training meaningfully improves cross-modal feature alignment and would manifest as a larger fusion gain on EAV. Concretely, AMCL would cluster intra-modal class representations, EMCL would align same-class representations across modalities, and SMCL would close the modality-specific embedding gap with a margin term $\alpha$. All three losses operate on the pre-classifier embeddings already exposed by the \texttt{forward()} contract, requiring no architectural changes.

\textbf{Phase 2.7: Identity-invariant representations without enrolment.} Subject calibration removes identity at inference time and requires an enrolment session. Adversarial identity confusion or identity-invariant contrastive objectives could remove it at training time instead, allowing one-shot inference on new users. The identity probe of Section~\ref{sec:cal_mechanism} provides the measurement, and the calibrated cross-subject results provide the target.

\textbf{Closing remark.} The thesis began by treating disagreement between the voluntary and involuntary modalities as the signal of interest, and ends with a more basic finding about what the signals carry. Every modality on EAV encodes who the participant is far more reliably than what they feel, and a fusion module trained on such features learns people as much as emotions. Removing the person from the representation, with nothing more than a short unlabelled recording, is what allows attention-based fusion to outperform averaging on people it has never seen.
